% ========================================================================
% فصل ۱۰: هوش مصنوعی و تشخیص ناهنجاری
% (بدنهٔ فصل؛ پشتیبانی همهٔ ادعاها از یادداشت‌های sam-work/source-review/)
% ========================================================================
\section{هوش مصنوعی و تشخیص ناهنجاری}

روش‌های هوش مصنوعی و یادگیری ماشین نقش مهمی در تقویت امنیت شبکهٔ اقتضایی خودرویی دارند؛ در مرور فارسیمدان و همکاران، علاقهٔ روزافزونی به به‌کارگیری یادگیری ماشین در امنیت شبکه‌های خودرویی برای تشخیص و پیش‌بینی بهتر، سریع‌تر و دقیق‌تر حملات پدید آمده است \cite{farsimadanReviewSecurityChallenges2025}. در طبقه‌بندی راهکارهای امنیتی این مرور (جدول ۱۰)، راهکارها به رمزنگاری و بلاکچین، یادگیری ماشین و عمیق، اعتماد، هویت، و \lr{5G/6G} دسته‌بندی می‌شوند و در شاخهٔ یادگیری ماشین نیز سه گروه نظارت‌شده، بدون نظارت و تقویتی طرح شده است \cite{farsimadanReviewSecurityChallenges2025}. این رویکردها با دوران کلان‌داده پیوند می‌خورند؛ اینترنت خودروها از اکتساب، انتقال، ذخیره‌سازی و پردازش کلان‌داده پشتیبانی می‌کند و در مقابل، کلان‌داده اینترنت خودروها را از نظر مشخصه‌سازی شبکه، تحلیل کارایی و طراحی پروتکل تقویت می‌کند \cite{xuInternetVehiclesBig2018}. این فصل این روش‌ها را در چهار حوزهٔ اصلی زیر بررسی می‌کند.

\subsection{تشخیص ناهنجاری و حملات}

نخستین و پرکاربردترین این حوزه‌ها، تشخیص ناهنجاری و حملات است. روش‌های کلاسیک تشخیص ناهنجاری در جریان داده بر سه رویکرد استوارند: یادگیری مبتنی بر قاعده که هر ارزیابی را بر پایهٔ دانش موجود اعتبارسنجی می‌کند، بررسی متقابل معیارهایی مانند سرعت و موقعیت میان جریان‌های دادهٔ گوناگون حسگرها، و مشاهدهٔ جریان داده در یک پنجرهٔ زمانی لغزان برای یافتن روندهای غیرعادی؛ با این حال، به دلیل حجم و تنوع زیاد اطلاعات، ساختن قواعد دستی اغلب دشوار است \cite{farsimadanReviewSecurityChallenges2025}.

یادگیری ماشین این مشکل را با رویکرد داده‌محور حل می‌کند: سامانهٔ تشخیص مبتنی بر یادگیری ماشین قواعد را خود از داده‌ها استخراج می‌کند، یک روش تشخیص واحد را برای سامانه‌های مختلف قابل بهره‌گیری می‌کند و توسعهٔ نرم‌افزار امنیتی را برای هر زیرسامانه ساده‌تر می‌سازد. در همین راستا، مدل‌های پیش‌خور و بازگشتی یادگیری عمیق برای تشخیص نفوذ به کار رفته‌اند؛ از جمله شبکهٔ عصبی عمیق برای شبکهٔ درون‌خودرویی و حافظهٔ کوتاه‌مدت ماندگار \LTRfootnote{\lr{Long Short-Term Memory (LSTM)}} برای تشخیص ناهنجاری داده‌های شبکهٔ کنترل خودرو \cite{farsimadanReviewSecurityChallenges2025}. مدل‌های یادگیری عمیق متشکل از لایه‌های شبکهٔ عصبی پیچشی \LTRfootnote{\lr{Convolutional Neural Network (CNN)}} و \lr{LSTM} می‌توانند با بهره‌گیری از داده‌های جریان‌های متمایز از نظر زمانی و جغرافیایی، تشخیصی مقاوم انجام دهند \cite{farsimadanReviewSecurityChallenges2025}. مرورهای جامعی نیز بر سامانه‌های تشخیص رفتار نادرست مبتنی بر یادگیری ماشین انجام شده است \cite{farsimadanReviewSecurityChallenges2025}.

در گروه یادگیری بدون نظارت، مدلی برای تشخیص پارازیت فرکانس رادیویی \lr{(radio frequency jamming)} در ارتباطات خودرو با همه‌چیز پیشنهاد شده است \cite{farsimadanReviewSecurityChallenges2025} و در کار دیگری، با یک راهبرد داده‌محور، حملهٔ منع سرویس و سه نوع حملهٔ دیگر علیه شبکهٔ درون‌خودرویی تشخیص داده می‌شود؛ در این روش ویژگی‌های مرتبط را با یادگیری بدون نظارت از اطلاعات گذرگاه شبکهٔ کنترل‌کننده \LTRfootnote{\lr{Controller Area Network (CAN)}} استخراج می‌کنند \cite{farsimadanReviewSecurityChallenges2025}. در کار دیگری نیز یادگیری ماشین برای تشخیص چهار نوع حمله به کنترل سرعت تطبیقی، یعنی حملات مبتنی بر سرعت، شتاب، موقعیت و ترکیب سرعت و موقعیت، به کار رفته است \cite{farsimadanReviewSecurityChallenges2025}.

فراتر از تشخیص نفوذ، مفهوم آگاهی موقعیتی \LTRfootnote{\lr{Situational Awareness}} در امنیت خودرویی مطرح است؛ این مفهوم شامل درک، ادراک و پیش‌بینی پویایی محیط است و در امنیت خودرویی به راهبردی امنیتی در ابر خودرویی اشاره دارد که معماری و کانال‌های شبکه را می‌شناسد. این رویکرد چالش‌برانگیزتر و پیچیده‌تر از تشخیص نفوذ است، زیرا باید در محیطی پویا با داده‌های ناپایدار پیش‌بینی کند \cite{farsimadanReviewSecurityChallenges2025}. افزون بر این، حتی بهترین راهبردهای دفاعی نیز در برابر شکست آسیب‌پذیرند؛ در نمونهٔ حملهٔ جعل \lr{GPS} موفق که خودرو را به سمت لبهٔ جاده هدایت می‌کند، خودرو باید بخش‌های امن و خطرناک جاده را تشخیص دهد و بخش‌بندی بلادرنگ صحنهٔ رانندگی با \lr{CNN} این سازوکار پشتیان (جایگزین) را فراهم می‌کند \cite{farsimadanReviewSecurityChallenges2025}. در عین حال، سامانه‌های مولدی مانند شبکه‌های مولد تخاصمی \LTRfootnote{\lr{Generative Adversarial Networks (GANs)}} می‌توانند برای ارزیابی سامانه‌های تشخیص، سیگنال‌های نفوذی آماریاً یکسان تولید کنند \cite{farsimadanReviewSecurityChallenges2025}.

\subsection{تحلیل آسیب‌پذیری و اعتماد}

تشخیص ناهنجاری تنها بخشی از امنیت مبتنی بر یادگیری است؛ در کنار آن، تحلیل آسیب‌پذیری و مدیریت اعتماد نقش مکملی دارند. بررسی بسته‌های ترافیک زندهٔ شبکه نیز برای شناسایی آسیب‌پذیری‌ها به کار رفته است \cite{farsimadanReviewSecurityChallenges2025}. اعتماد به اطمینان یک موجودیت به موجودیت دیگری در شبکه گفته می‌شود و محاسبهٔ آن معیار امنیتی تکمیلی است \cite{farsimadanReviewSecurityChallenges2025}.

در این حوزه، روشی برای محاسبهٔ اعتماد مبتنی بر منطق فازی \LTRfootnote{\lr{Fuzzy Logic}} پیشنهاد شده است که دقت و یکپارچگی پیام‌های رویداد و فرستندگان آن‌ها را ارزیابی می‌کند. در این رویکرد ابتدا مجموعه‌ها، الزامات و معیارهای فازی تعریف می‌شوند، سپس مقادیر متغیرهای ورودی مقداردهی اولیه می‌شوند و در پایان سامانهٔ فازی با به‌کارگیری قواعد فازی، خروجی را مشخص و نتایج را ارزیابی می‌کند \cite{farsimadanReviewSecurityChallenges2025}.

\subsection{یادگیری توزیع‌شده و تقویتی}

حوزهٔ بعدی، یادگیری توزیع‌شده و تقویتی است که آموزش مدل و بهینه‌سازی سیاست امنیتی را فراتر از یک گره می‌برد. در بافت حریم خصوصی، مدل یادگیری فدرال \LTRfootnote{\lr{Federated Learning (FL)}} مبتنی بر بلاکچین برای انتقال امن اطلاعات، با توجه به نگرانی ارائه‌دهندگان دربارهٔ حریم خصوصی، پیشنهاد شده است؛ لو و همکاران این مدل را به‌صورت ناهمگام طراحی کرده‌اند \cite{farsimadanReviewSecurityChallenges2025}. در رویکردی دیگر، خودروها تجربه را برای بهبود تشخیص خودروهای رفتار نادرست با یکدیگر به اشتراک می‌گذارند؛ البته همکاری میان خودروها نگرانی‌های حریم خصوصی نیز ایجاد می‌کند \cite{farsimadanReviewSecurityChallenges2025}.

یادگیری تقویتی \LTRfootnote{\lr{Reinforcement Learning (RL)}} نیز کاربردهای متعددی در امنیت خودرویی دارد. در شبکه‌های خودرویی نرم‌افزارمحور، این روش برای تعیین امن‌ترین سیاست ارتباط با توجه به اثر خودروهای بدخواه به کار رفته است \cite{farsimadanReviewSecurityChallenges2025}. رویکرد یادگیری \lr{Q} نیز برای ارزیابی اعتمادپذیری خودروهای خودران و شناسایی نفوذگران بر پایهٔ سطح اعتماد محاسبه‌شده، با دو رویکرد ارزیابی مستقیم و غیرمستقیم، پیشنهاد شده است \cite{farsimadanReviewSecurityChallenges2025}. در سازوکاری محور داده نیز خودرو از مدل ارزیابی اعتماد، مقدار اعتماد رویدادهایی چون تنظیم سرعت و انتخاب مسیر را می‌پرسد \cite{farsimadanReviewSecurityChallenges2025} و سامانهٔ احراز اصالت فیزیکی مبتنی بر یادگیری تقویتی نیز برای محافظت در برابر حملات جعل ناشناس پیشنهاد شده است \cite{farsimadanReviewSecurityChallenges2025}.

ترکیب هوش مصنوعی با بلاکچین اینترنت خودروها را غیرمتمرکز، هوشمند و امن می‌کند؛ با این حال، هوش مصنوعی در برابر حملات سایبری آسیب‌پذیر است و بلاکچین سربار محاسباتی را افزایش می‌دهد، بنابراین معماری‌ای پیشنهاد می‌شود که هوش مصنوعی و بلاکچین را در هم بیامیزد و از رایانش لبهٔ خودرویی بهره بگیرد \cite{hozouriOverviewVANETVehicular2023}. از سوی دیگر، پیش‌بینی جریان ترافیک با یادگیری عمیق در ابر خودرویی گزارش شده است \cite{farsimadanReviewSecurityChallenges2025} و کلان‌داده‌های خودرویی برای پیش‌بینی جریان ترافیک و مسیر کوتاه‌مدت نیز مورد تحلیل قرار می‌گیرند \cite{xuInternetVehiclesBig2018}.

\subsection{چالش‌ها}

با وجود این پیشرفت‌ها، چالش‌هایی در داده، ارزیابی و پیاده‌سازی باقی است. روش‌های مبتنی بر یادگیری فدرال، هرچند پیشرفت‌هایی در سامانه‌های حفظ‌کنندهٔ حریم خصوصی ارتباط خودرو با همه‌چیز دارند، همچنان به ساختار کارگزار و کارخواه وابسته‌اند که آن‌ها را در برابر فعالیت‌های مخرب آسیب‌پذیر می‌کند؛ علاوه بر این، راهبردهای تخاصمی و مسموم‌سازی داده ممکن است بر مدل‌هایی که به‌صورت محلی آموزش دیده‌اند اثر بگذارند و طراحی مدل‌های مقاوم و حفظ‌کنندهٔ حریم خصوصی مبتنی بر یادگیری فدرال از جهت‌های آیندهٔ پژوهش در این حوزه است \cite{farsimadanReviewSecurityChallenges2025}. حفظ حریم خصوصی در یادگیری‌های مشارکتی نیز چالش‌برانگیز است؛ همکاری میان خودروها برای بهبود تشخیص، نگرانی‌های حریم خصوصی ایجاد می‌کند \cite{farsimadanReviewSecurityChallenges2025}.

ارزیابی این روش‌ها نیز واقع‌گرایی کافی ندارد؛ بیشتر آن‌ها با شبیه‌سازی ارزیابی شده‌اند و شبیه‌سازها ضعیف‌اند، و بیشتر پژوهشگران داده‌های مورد استفادهٔ خود را در ارزیابی‌ها به‌صورت عمومی منتشر نمی‌کنند؛ بنابراین نیاز به مجموعه‌داده‌های باز برای پیشبرد پژوهش‌های یادگیری‌محور در این حوزه همچنان احساس می‌شود \cite{farsimadanReviewSecurityChallenges2025}. افزون بر این، توان محاسباتی محدود خودروها یک چالش پایه‌ای است و اجرای برخی راهکارهای امنیتی خودرویی به دلیل سربار بالای آن‌ها ممکن نیست \cite{farsimadanReviewSecurityChallenges2025}؛ سربار محاسباتی بالای بلاکچین نیز در ترکیب با هوش مصنوعی باید در طراحی معماری در نظر گرفته شود \cite{hozouriOverviewVANETVehicular2023}.
